\section{Yêu cầu và cơ sở lý thuyết}

\subsection{Các đại lượng giám sát IAQ}

Trong phạm vi đề tài, hệ thống tập trung theo dõi năm nhóm đại lượng chính gồm
CO\textsubscript{2}, bụi mịn, VOC, nhiệt độ và độ ẩm. Các đại lượng này được lựa chọn
vì phản ánh đồng thời tình trạng thông gió, mức độ ô nhiễm không khí và điều kiện tiện
nghi trong phòng học.

CO\textsubscript{2} được sử dụng như một chỉ báo gián tiếp của mức độ thông gió và mật
độ người trong phòng. Khi lớp học đông người hoặc thông gió kém, nồng độ
CO\textsubscript{2} có xu hướng tăng dần, do đó việc theo dõi đại lượng này giúp hệ
thống nhận biết thời điểm cần tăng thông gió hoặc phát cảnh báo. Đây cũng là chỉ báo
được sử dụng phổ biến trong các tài liệu hướng dẫn IAQ cho trường học
\cite{epaiaqschools,alaiaqschools}.

Bụi mịn được biểu diễn qua các chỉ số PM1.0, PM2.5 và PM10. Trong hệ thống này, cảm
biến PMS7003 có khả năng đọc cả ba chỉ số, tuy nhiên pipeline IoT hiện tại ưu tiên sử
dụng PM2.5 làm chỉ số bụi chính để publish MQTT, lưu cơ sở dữ liệu và hiển thị trên
dashboard. PM2.5 được lựa chọn vì liên quan trực tiếp đến nguy cơ sức khỏe và thường
được dùng trong các khuyến nghị chất lượng không khí \cite{whoaqg2021}.

VOC phản ánh sự hiện diện của các hợp chất hữu cơ bay hơi phát sinh từ vật liệu xây
dựng, sơn, dung môi, hóa chất vệ sinh hoặc hoạt động sinh hoạt trong phòng. Một số hợp
chất thuộc nhóm này đã được WHO xếp vào nhóm chất ô nhiễm trong nhà cần quan tâm
\cite{whoiaq}. Trong nguyên mẫu hiện tại, giá trị VOC được đọc ở dạng \textit{raw/index}
tương đối, do đó chủ yếu được dùng để theo dõi xu hướng và kích hoạt logic cảnh báo ở
mức prototype.

Nhiệt độ và độ ẩm không trực tiếp đại diện cho nồng độ chất ô nhiễm, nhưng ảnh hưởng
rõ rệt đến tiện nghi nhiệt, cảm nhận môi trường và điều kiện vận hành của lớp học. Vì
vậy, hai đại lượng này được đưa vào hệ thống nhằm hỗ trợ đánh giá tổng thể điều kiện
môi trường và làm cơ sở cho các quyết định điều khiển cục bộ.

\subsection{Định hướng giám sát và cảnh báo}

Mục tiêu của hệ thống không phải là thay thế thiết bị đo kiểm chuyên dụng hay đưa ra kết
luận y tế tuyệt đối về chất lượng không khí. Đề tài hướng đến việc xây dựng một nguyên
mẫu IoT có khả năng giám sát liên tục, phát hiện xu hướng bất thường và hỗ trợ điều
khiển cục bộ. Vì vậy, dữ liệu đo được được sử dụng theo hai hướng chính: hiển thị trạng
thái môi trường theo thời gian thực và kích hoạt cảnh báo hoặc thiết bị chấp hành khi
đại lượng vượt khỏi vùng vận hành mong muốn.

Đối với CO\textsubscript{2}, PM2.5, nhiệt độ và độ ẩm, hệ thống có thể so sánh trực
tiếp với các ngưỡng cấu hình trong firmware. Đối với VOC, do giá trị hiện tại mang tính
tương đối, hệ thống chủ yếu dùng để quan sát xu hướng và kích hoạt cảnh báo ở mức
nguyên mẫu. Cách tiếp cận này phù hợp với phạm vi đồ án, vốn tập trung vào kiến trúc đo,
truyền dữ liệu, xử lý lỗi, hiển thị và điều khiển hơn là chứng nhận IAQ theo tiêu chuẩn
y tế hoặc công trình.

Trong firmware RTOS, các cặp ngưỡng bật/tắt như \texttt{co2\_on/co2\_off},
\texttt{pm\_on/pm\_off}, \texttt{voc\_on/voc\_off}, \texttt{temp\_high\_on/temp\_high\_off}
và các ngưỡng độ ẩm thấp/cao được dùng cùng cơ chế hysteresis để tránh thiết bị bật tắt
liên tục khi giá trị đo dao động gần ngưỡng. Các tham số này có thể được hiệu chỉnh
trong quá trình kiểm thử thông qua backend hoặc giao diện điều khiển.

\subsection{Yêu cầu chức năng}

Từ mục tiêu giám sát IAQ và điều khiển thiết bị trong phòng học, hệ thống cần đáp ứng
các yêu cầu chức năng sau:

\begin{enumerate}[leftmargin=2em]
    \item ESP32 đọc dữ liệu cảm biến theo chu kỳ và tạo mẫu đo có dấu thời gian tương
    đối trong quá trình vận hành.
    \item Firmware thu thập được dữ liệu từ nhiều cảm biến sử dụng các giao tiếp khác
    nhau, bao gồm cảm biến bụi, CO\textsubscript{2}, VOC, nhiệt độ và độ ẩm.
    \item Dữ liệu cảm biến được chuẩn hóa về một mô hình telemetry thống nhất trước khi
    đưa vào các khối xử lý, điều khiển hoặc truyền thông MQTT.
    \item Firmware biểu diễn được trạng thái dữ liệu thiếu, lỗi đọc hoặc mất kết nối để
    tránh sử dụng giá trị không đáng tin cậy cho quyết định điều khiển.
    \item Hệ thống đánh giá từng nhóm đại lượng IAQ dựa trên ngưỡng cấu hình và xác
    định nguyên nhân gây cảnh báo.
    \item Hệ thống điều khiển thiết bị cục bộ theo nguyên nhân phát hiện được, chẳng
    hạn thông gió cho CO\textsubscript{2}/VOC, lọc bụi cho PM2.5, điều hòa cho nhiệt
    độ hoặc tạo ẩm khi độ ẩm thấp.
    \item ESP32 gửi telemetry lên backend thông qua MQTT theo định dạng JSON thống nhất.
    \item Backend tiếp nhận telemetry, kiểm tra định dạng dữ liệu, lưu lịch sử vào cơ sở
    dữ liệu và cung cấp API cho dashboard.
    \item Dashboard web hiển thị dữ liệu realtime, biểu đồ lịch sử, trạng thái cảm biến,
    trạng thái điều khiển tự động và các cảnh báo đang kích hoạt.
    \item Dashboard cho phép bật/tắt chế độ tự động, điều khiển thủ công từng thiết bị
    khi ở manual mode và cập nhật ngưỡng điều khiển của hệ thống.
\end{enumerate}

\subsection{Yêu cầu phi chức năng}

Bên cạnh các chức năng chính, hệ thống cần đáp ứng một số yêu cầu phi chức năng để vận
hành ổn định trên nền tảng nhúng và phù hợp với mô hình IoT.

Thứ nhất, firmware cần nhẹ và ổn định trên ESP32. Các tác vụ đọc cảm biến, xử lý dữ
liệu, điều khiển thiết bị và truyền thông không được làm treo toàn bộ hệ thống khi một
cảm biến mất kết nối hoặc phản hồi chậm. Vì vậy, driver cần có timeout, trạng thái lỗi
rõ ràng và cơ chế bỏ qua mẫu đo không hợp lệ.

Thứ hai, dữ liệu nội bộ cần được biểu diễn nhất quán và tiết kiệm tài nguyên. Trong
firmware hiện tại, một số đại lượng được xử lý nội bộ dưới dạng fixed-point nhân hệ số
10 trước khi chuyển về giá trị phù hợp để publish hoặc hiển thị.

Thứ ba, hệ thống cần có khả năng mở rộng. Việc bổ sung cảm biến mới, thiết bị chấp hành
mới hoặc chỉ số IAQ mới không nên làm thay đổi toàn bộ kiến trúc firmware và backend,
vì vậy các khối đọc cảm biến, chuẩn hóa mẫu đo, đánh giá IAQ, truyền MQTT và hiển thị
dashboard cần được tách tương đối độc lập.

Thứ tư, tầng truyền thông IoT cần có định dạng dữ liệu rõ ràng và dễ kiểm thử. MQTT
được sử dụng cho luồng publish/subscribe giữa ESP32 và backend vì đây là giao thức nhẹ,
phù hợp với các hệ thống IoT và thiết bị có tài nguyên hạn chế~\cite{oasismqtt}. Payload
JSON cần có cấu trúc ổn định, có mã thiết bị, thời gian, giá trị cảm biến và trạng thái
điều khiển để backend có thể kiểm tra trước khi lưu trữ.

Thứ năm, backend và dashboard cần hỗ trợ truy vết dữ liệu theo thời gian. Dữ liệu lịch
sử cần được lưu kèm timestamp để phục vụ biểu đồ, đối chiếu khi kiểm thử và phân tích
xu hướng IAQ. Đồng thời, các cấu hình nhạy cảm như địa chỉ broker hoặc thông tin cơ sở
dữ liệu cần được tách khỏi mã nguồn thông qua biến môi trường hoặc tệp cấu hình cục bộ.

Thứ sáu, dashboard cần thể hiện rõ trạng thái cảm biến, trạng thái dữ liệu, trạng thái
thiết bị và chế độ điều khiển để người dùng có thể quan sát và thao tác nhất quán.
